\subsubsection{Proof strategy towards Theorem \ref{thm:b0ConcGen}} 
%
The proof of Theorem \ref{thm:b0ConcGen} is the main goal of the paper. 
%
We will assume for the sake of contradiction the existence of a (small) parameter $\delta \in (0,1)$ and an infinite increasing sequence $(x_j)_j$ such that
\begin{equation}\label{eq:converseCond}
\frac{1}{\log x_j} \sum_{\ss{n \leq x_j \\ f_1(an+b)f_2(cn+d) \neq 0}} \frac{1_{f_1(an+b) = Cf_2(cn+d)}}{n} \geq \delta, \quad j \geq 1.
\end{equation}
%
%
We seek a contradiction by obtaining from \eqref{eq:converseCond} precise information about the sequence $(|f(p)|)_p$. 

In order to analyse the condition \eqref{eq:converseCond} we 
%
appeal to a deep theorem of Tao \cite{Tao} (see Theorem \ref{thm1:Tao} below) on logarithmically-averaged correlations of multiplicative functions $g_1,g_2$ taking values in $\mb{U}$, the closed unit disc. This result implies structural information about the prime-indexed sequences $(g_j(p))_{p}$, $j = 1,2$, whenever 
%
$$
\lim_{x \ra \infty} \frac{1}{\log x} 
\left|
\sum_{n \leq x} \frac{g_1(an+b)g_2(cn+d)}{n}
\right| \neq 0.
$$
%
To apply this we ``project'' our multiplicative functions $f_1,f_2$ onto $\mb{U}$ in a manner that retains their multiplicativity. 
%
%
%
%
%

For ease of exposition, 
%
suppose that our affine linear forms are $(an+b,cn+d) = (n,n+h)$, $h \geq 1$, that $C = 1$ and that $f_1 = f_2 = f$. 
%
We observe that for $t \in \mb{R}$ 
%
the composition
$$
|f|_t(n) := \begin{cases} |f(n)|^{it} &\text{ if $f(n) \neq 0$} \\ 0 &\text{ if $f(n) = 0$} \end{cases}
%
$$
is a well-defined multiplicative function that takes values in $\mb{U}$. \\
%
%
%
%
Let $A \geq 1$ be a parameter to be chosen later. The function
%
$$
g(n) := 
%
A\log |f(n)| 
%
$$
is well-defined on the set $\{n \in \mb{N} : f(n) \neq 0\}$, and is additive when restricted to that set. In particular, we have
$$
f(n) = f(n+h) \neq 0 \Rightarrow g(n) = g(n+h),
$$
and also\footnote{Here we use the standard notation $e(y) := e^{2\pi i y}$ for $y \in \mb{R}$.} $|f|_{2\pi At}(n) = e(t g(n))$. Using a Fourier analytic identity, we thus show that if \eqref{eq:converseCond} holds then there is a bounded set $X \subset \mb{R}$ with positive Lebesgue measure such that for each $y \in X$ we have lower bounds 
$$
\frac{1}{\log x_k} \left|\sum_{n \leq x_k} \frac{|f|_{2\pi A y}(n) |f|_{-2\pi A y}(n+h)}{n}\right| \gg_{\delta} 1.
$$
